\chapter{Circuitos RC e RL}\label{ch:g12:rc-rl-circuits}

Aperte o disparador de uma câmera até a metade: um zunido fino sobe por dois segundos,
uma luz fica verde --- e só então o flash dispara, despejando em um milissegundo o que
levou dois segundos juntando. Este capítulo apresenta os dois componentes que dão aos
circuitos uma memória e um relógio --- o \omterm{def:g12:rc-rl-circuits:capacitor}{capacitor} e a \omterm{def:g11:magnetic-fields:solenoid}{bobina} --- e as primeiras
equações diferenciais da física, conferíveis com as derivadas deste ano.

\section{O capacitor}

\begin{definition}[Capacitor e capacitância]\label{def:g12:rc-rl-circuits:capacitor}
Um \emph{capacitor}\index{capacitor} é um par de placas metálicas frente a frente,
separadas por um isolante fino. Ligadas a uma fonte, as placas recebem cargas opostas
$+q$ e $-q$ --- um \omterm{def:g11:electric-gravitational-fields:efield}{campo elétrico} preenche o vão
(\cref{ch:g11:electric-gravitational-fields}) --- e a tensão $u$ entre as placas é
proporcional à carga armazenada:
\[ q = C\,u . \]
A constante $C$ é a \emph{capacitância}\index{capacitância}, medida em
\emph{farads}\index{farad} (\unit{F}): $\qty{1}{F} = \qty{1}{C/V}$.
\end{definition}

\begin{example}[Ordens de grandeza]\label{ex:g12:rc-rl-circuits:orders}
Um \omterm{def:g12:rc-rl-circuits:capacitor}{farad} é enorme --- um \omterm{def:g11:fundamental-interactions:charge}{coulomb} estacionado por \omterm{def:g11:circuits-and-power:voltage}{volt}. Os cerâmicos dos circuitos de
rádio: de \unit{pF} a \unit{nF}; os eletrolíticos das fontes de alimentação e dos
flashes: de \unit{\micro F} a \unit{mF}; os supercapacitores: milhares de \omterm{def:g12:rc-rl-circuits:capacitor}{farads}.
\end{example}

\begin{proposition}[Corrente que entra num capacitor]\label{prop:g12:rc-rl-circuits:cdudt}
A corrente que chega à placa de um \omterm{def:g12:rc-rl-circuits:capacitor}{capacitor} é a taxa com que a carga dele cresce:
\[
i = \frac{dq}{dt} = C\,\frac{du}{dt} ,
\]
com as minúsculas marcando as grandezas que variam no tempo ($U$ e $I$ permanecem
constantes, \cref{ch:g11:circuits-and-power}).
\end{proposition}

\begin{proof}
Corrente é carga entregue por segundo (\cref{def:g11:circuits-and-power:current}): num
intervalo que encolhe, $i$ é a derivada de $q$; e $q = Cu$ com $C$ constante dá
$C\,du/dt$.
\end{proof}

\begin{remark}[Duas consequências]\label{rem:g12:rc-rl-circuits:continuity}
Se $u$ é constante, $i = 0$: um \omterm{def:g12:rc-rl-circuits:capacitor}{capacitor} carregado não deixa passar corrente contínua.
E $u$ nunca pode dar um salto --- um degrau exigiria corrente infinita: a tensão do
\omterm{def:g12:rc-rl-circuits:capacitor}{capacitor} é a memória do circuito.
\end{remark}

\section{Carregar um capacitor: o circuito RC}

Uma única malha --- fonte $\mathcal{E}$ (\cref{def:g11:circuits-and-power:emf}), chave,
resistor, \omterm{def:g12:rc-rl-circuits:capacitor}{capacitor} vazio: feche $K$ em $t = 0$.

\begin{omfigure}
\begin{circuitikz}[font=\small]
  \draw (0,0) to[battery1, l=$\mathcal{E}$] (0,2.4) to[nos, l=$K$] (2.1,2.4)
    to[R=$R$, i=$i$] (4.2,2.4) -- (5.4,2.4) to[C=$C$, v=$u$] (5.4,0) -- (0,0);
\end{circuitikz}

{\small O \omterm{thm:g12:rc-rl-circuits:charging}{circuito RC} de carga: fechar $K$ deixa a fonte empurrar carga através de $R$
até as placas de $C$.}
\end{omfigure}

\begin{theorem}[Lei de carga do circuito RC]\label{thm:g12:rc-rl-circuits:charging}
\index{circuito RC}
Depois que a chave fecha, a tensão do \omterm{def:g12:rc-rl-circuits:capacitor}{capacitor} obedece à equação abaixo e, a partir de
$u(0) = 0$, decorre dela:
\[
RC\,\frac{du}{dt} + u = \mathcal{E},
\qquad
u(t) = \mathcal{E}\left(1 - e^{-t/\tau}\right),
\qquad \tau = RC .
\]
\end{theorem}

\begin{proof}
As tensões ao longo da malha se somam (\cref{prop:g11:circuits-and-power:laws}):
$\mathcal{E} = u_R + u = Ri + u$ a cada instante; substitua $i = C\,du/dt$
(\cref{prop:g12:rc-rl-circuits:cdudt}). Agora \emph{verifique} a candidata:
$du/dt = (\mathcal{E}/\tau)e^{-t/\tau}$, logo
$RC\,(\mathcal{E}/RC)\,e^{-t/\tau} + \mathcal{E}(1 - e^{-t/\tau}) = \mathcal{E}$,
verdadeiro para todo $t$; e $u(0) = \mathcal{E}(1-1) = 0$. Que nenhuma \emph{outra}
curva se ajusta à equação e à partida é demonstrado no curso de matemática deste ano,
sobre equações diferenciais.
\end{proof}

\begin{definition}[Constante de tempo]\label{def:g12:rc-rl-circuits:tau}
A \emph{constante de tempo}\index{constante de tempo} de um \omterm{thm:g12:rc-rl-circuits:charging}{circuito RC} é $\tau = RC$.
\omterm{def:g10:orders-of-magnitude:unit}{Unidades}: $\unit{\ohm} \times \unit{F} = (\unit{V/A})(\unit{C/V}) = \unit{C/A} =
\unit{s}$ --- o compasso próprio do circuito.
\end{definition}

\begin{method}[Ler $\tau$ na curva]\label{met:g12:rc-rl-circuits:reading}
Três leituras equivalentes, direto de um gráfico:
\begin{enumerate}
  \item \emph{tangente na origem:} $u'(0) = \mathcal{E}/\tau$, de modo que a tangente
        inicial alcança a assíntota em $t = \tau$;
  \item \emph{63\%:} $u(\tau) = \mathcal{E}(1 - e^{-1}) \approx 0.63\,\mathcal{E}$;
  \item \emph{95\%:} $u(3\tau) \approx 0.95\,\mathcal{E}$ --- carregado, para fins
        práticos ($5\tau$: $99\%$).
\end{enumerate}
\end{method}

\begin{proposition}[Descarga]\label{prop:g12:rc-rl-circuits:discharge}
Um \omterm{def:g12:rc-rl-circuits:capacitor}{capacitor} carregado a $U_0$ e fechado sobre um resistor $R$ sozinho obedece a
$RC\,du/dt + u = 0$ e segue
\[
u(t) = U_0\,e^{-t/\tau}, \qquad \tau = RC
\]
--- restam $37\%$ em $\tau$ e $5\%$ em $3\tau$.
\end{proposition}

\begin{proof}
A mesma lei das malhas, sem fonte; e de novo a substituição:
$RC(-U_0/\tau)e^{-t/\tau} + U_0 e^{-t/\tau} = 0$, com $u(0) = U_0$.
\end{proof}

\begin{omfigure}
\begin{tikzpicture}
\begin{axis}[omaxis, width=9.2cm, height=5.4cm, xlabel={$t$}, ylabel={$u$},
    xmin=0, xmax=5, ymin=0, ymax=7.4, xtick={0,1,2,3,4,5},
    xticklabels={$0$,$\tau$,$2\tau$,$3\tau$,$4\tau$,$5\tau$},
    ytick={0,2.21,3.79,6.0},
    yticklabels={$0$,$0.37\,\mathcal{E}$,$0.63\,\mathcal{E}$,$\mathcal{E}$}]
  \addplot[omExo, dashed, domain=0:5, samples=2] {6};
  \addplot[omExo, dashed, domain=0:1.18, samples=2] {6*x};
  \addplot[omExo, dashed, domain=0:1.1, samples=2] {6-6*x};
  \addplot[omDef, thick, domain=0:5, samples=80] {6*(1-exp(-x))};
  \addplot[omThm, thick, domain=0:5, samples=80] {6*exp(-x)};
  \draw[dashed, omExo] (1,0) -- (1,3.79) -- (0,3.79);
  \draw[dashed, omExo] (3,0) -- (3,5.70);
\end{axis}
\end{tikzpicture}

{\small Carga (azul, de $u = 0$ rumo a $\mathcal{E}$) e descarga (vermelho, a partir de
$U_0 = \mathcal{E}$). As duas tangentes iniciais encontram a assíntota delas em
$t = \tau$; $63\%$ feito em $\tau$, $95\%$ em $3\tau$.}
\end{omfigure}

\begin{example}[Os números de um flash]\label{ex:g12:rc-rl-circuits:flashrc}
Um flash fotográfico carrega $C = \qty{800}{\micro F}$ através de
$R = \qty{1.0}{k\ohm}$: $\tau = 10^3 \times 8.0\times10^{-4} = \qty{0.80}{s}$, de modo
que a luz de ``pronto'' --- ligada para acender a $95\%$ --- espera
$3\tau \approx \qty{2.4}{s}$: o zunido da espera. Disparado através dos meros
\qty{0.50}{\ohm} do tubo de flash, o \emph{mesmo} \omterm{def:g12:rc-rl-circuits:capacitor}{capacitor} se esvazia com
$\tau' = \qty{0.40}{ms}$: duas mil vezes mais rápido para sair do que para entrar.
\end{example}

\section{A bobina: indutância e o circuito RL}

\begin{definition}[Indutor e indutância]\label{def:g12:rc-rl-circuits:inductor}
Um \emph{indutor}\index{indutor} é uma \omterm{def:g11:magnetic-fields:solenoid}{bobina} de fio: a corrente dele enfia nas próprias
espiras um \omterm{def:g11:magnetic-fields:bfield}{campo magnético} (\cref{ch:g11:magnetic-fields}); sempre que esse \omterm{def:g11:electric-gravitational-fields:field}{campo}
\emph{muda}, aparece uma tensão nos terminais da \omterm{def:g11:magnetic-fields:solenoid}{bobina}:
\[
u = L\,\frac{di}{dt} .
\]
A constante $L$ é a \emph{indutância}\index{indutância}, medida em
\emph{henries}\index{henry} (\unit{H}): $\qty{1}{H} = \qty{1}{V.s/A}$.
\end{definition}
\admitted

\begin{remark}[A bobina resiste à mudança]\label{rem:g12:rc-rl-circuits:oppose}
Por que uma corrente que varia gera uma tensão --- a indução --- é deduzido nos volumes
universitários; aqui tomamos a lei e o caráter dela: a \omterm{def:g11:magnetic-fields:solenoid}{bobina} \emph{se opõe às
mudanças} da corrente dela. Gêmeos espelhados: o $u$ de um \omterm{def:g12:rc-rl-circuits:capacitor}{capacitor} não pode saltar, o
$i$ de uma \omterm{def:g11:magnetic-fields:solenoid}{bobina} não pode saltar.
\end{remark}

\begin{omfigure}
\begin{circuitikz}[font=\small]
  \draw (0,0) to[battery1, l=$\mathcal{E}$] (0,2.4) to[nos, l=$K$] (2.1,2.4)
    to[R=$R$] (4.2,2.4) -- (5.4,2.4) to[L=$L$, i=$i$, v=$u$] (5.4,0) -- (0,0);
\end{circuitikz}

{\small O \omterm{prop:g12:rc-rl-circuits:rl}{circuito RL}: depois que $K$ fecha, a \omterm{def:g11:magnetic-fields:solenoid}{bobina} deixa a corrente subir apenas no
compasso próprio dela, $\tau = L/R$.}
\end{omfigure}

\begin{proposition}[Subida da corrente num circuito RL]\label{prop:g12:rc-rl-circuits:rl}
\index{circuito RL}
Fechar a chave sobre uma fonte $\mathcal{E}$, um resistor $R$ e uma \omterm{def:g11:magnetic-fields:solenoid}{bobina} $L$ em série
dá $\mathcal{E} = Ri + L\,di/dt$, ou seja,
\[
\frac{L}{R}\,\frac{di}{dt} + i = \frac{\mathcal{E}}{R},
\qquad\text{resolvida por}\quad
i(t) = \frac{\mathcal{E}}{R}\left(1 - e^{-t/\tau}\right),
\quad \tau = \frac{L}{R} .
\]
\end{proposition}

\begin{proof}
De novo a lei das malhas; e é a \emph{mesma equação} do
\cref{thm:g12:rc-rl-circuits:charging} com $u \to i$,
$\mathcal{E} \to \mathcal{E}/R$ e $RC \to L/R$ --- a mesma substituição verifica a
solução, e toda leitura do \cref{met:g12:rc-rl-circuits:reading} se transfere.
\end{proof}

\begin{example}[Um eletroímã acorda]\label{ex:g12:rc-rl-circuits:rl}
A \omterm{def:g11:magnetic-fields:solenoid}{bobina} de um relé, $L = \qty{0.30}{H}$ e $R = \qty{6.0}{\ohm}$, numa fonte de
\qty{12}{V}: corrente final $\mathcal{E}/R = \qty{2.0}{A}$, no compasso
$\tau = 0.30/6.0 = \qty{50}{ms}$ --- desperto em $3\tau \approx \qty{0.15}{s}$.
\end{example}

\begin{remark}[A faísca na abertura]\label{rem:g12:rc-rl-circuits:spark}
Abra uma chave sobre uma \omterm{def:g11:magnetic-fields:solenoid}{bobina} energizada e $i$ precisa despencar a zero numa fração
de milissegundo: $di/dt$ fica enorme e $u = L\,di/dt$ alcança centenas ou milhares de
\omterm{def:g11:circuits-and-power:voltage}{volts} --- uma faísca salta o contato. A \omterm{def:g11:magnetic-fields:solenoid}{bobina} de ignição de um carro faz isso de
propósito, milhares de vezes por minuto.
\end{remark}

\section{Energia armazenada}

\begin{proposition}[Energia num capacitor, energia numa bobina]\label{prop:g12:rc-rl-circuits:energy}
\index{energia!de um capacitor}\index{energia!de uma bobina}
Um \omterm{def:g12:rc-rl-circuits:capacitor}{capacitor} $C$ carregado à tensão $u$ e uma \omterm{def:g11:magnetic-fields:solenoid}{bobina} $L$ percorrida por uma corrente
$i$ armazenam
\[
E_C = \tfrac12\,C u^2 = \frac{q^2}{2C},
\qquad
E_L = \tfrac12\,L i^2 .
\]
\end{proposition}

\begin{proof}
Pousar uma pequena carga $dq$ em placas já à tensão $u$ custa $u\,dq$
(\cref{def:g11:circuits-and-power:voltage}); empilhando os custos, o total é a área sob
a reta $u = q/C$ --- o triângulo abaixo --- logo $E_C = \tfrac12 qu = \tfrac12 Cu^2$. O
mesmo triângulo para a \omterm{def:g11:magnetic-fields:solenoid}{bobina} (elevar $i$ de $di$ custa $Li\,di$) dá
$E_L = \tfrac12 Li^2$.
\end{proof}

\begin{omfigure}
\begin{tikzpicture}[font=\small, x=1.35cm, y=1.05cm]
  \fill[omDef!25] (0,0) -- (4,0) -- (4,2.6) -- cycle;
  \fill[omDef!60] (2.4,0) rectangle (2.7,1.658);
  \draw[omDef, thick] (0,0) -- (4,2.6);
  \draw[-Stealth] (0,0) -- (4.8,0) node[right] {$q$};
  \draw[-Stealth] (0,0) -- (0,3.1) node[above] {$u = q/C$};
  \draw[dashed] (4,0) -- (4,2.6) -- (0,2.6);
  \node[below] at (4,0) {$q_0$}; \node[left] at (0,2.6) {$u_0$};
  \node[below, font=\footnotesize] at (2.55,0) {$dq$};
  \node[omDef!50!black, font=\footnotesize] at (2.9,0.85)
    {$E_C$};
\end{tikzpicture}

{\small Cada fatia $dq$ custa $u\,dq$; a carga inteira custa a área do triângulo,
$E_C = \tfrac12 q_0 u_0$ --- os primeiros \omterm{def:g11:fundamental-interactions:charge}{coulombs} embarcam com tensão baixa.}
\end{omfigure}

\begin{example}[Ordens de grandeza de energia]\label{ex:g12:rc-rl-circuits:energies}
O \omterm{def:g12:rc-rl-circuits:capacitor}{capacitor} de flash do \cref{ex:g12:rc-rl-circuits:flashrc}:
$\tfrac12 \times 8.0\times10^{-4} \times 300^2 = \qty{36}{J}$. Um supercapacitor de
\qty{3000}{F} a \qty{2.7}{V}: \qty{1.1e4}{J} --- o equivalente a uma pilha AA. Uma
\omterm{def:g11:magnetic-fields:solenoid}{bobina} de \qty{10}{mH} a \qty{10}{A}: \qty{0.50}{J}. Menos do que as pilhas guardam ---
mas liberável em um milissegundo.
\end{example}

\begin{remark}[Para que servem os dois marcadores de tempo]\label{rem:g12:rc-rl-circuits:applications}
\index{filtragem}\index{circuito temporizador}
Entrar devagar e sair depressa é o que move o flash da câmera e o desfibrilador. O
\omterm{def:g12:mechanical-waves:delay}{atraso} $\tau = RC$ toca os \omterm{rem:g12:rc-rl-circuits:applications}{circuitos temporizadores}: limpadores de para-brisa
intermitentes, balizas piscantes, luzes de escada. Um \omterm{def:g12:rc-rl-circuits:capacitor}{capacitor} em paralelo com uma
fonte \emph{filtra} a tensão ($u$ não pode saltar); uma \omterm{def:g11:magnetic-fields:solenoid}{bobina} em série filtra a
corrente ($i$ não pode saltar). E \omterm{def:g11:magnetic-fields:solenoid}{bobina} mais \omterm{def:g12:rc-rl-circuits:capacitor}{capacitor} juntos fazem um \omterm{def:g12:oscillators-and-time:oscillator}{oscilador}
(\cref{ch:g12:rlc-oscillations}).
\end{remark}

\section{Exercícios}

\begin{exercise}[$\star$]\label{exo:g12:rc-rl-circuits:1}
Um \omterm{def:g12:rc-rl-circuits:capacitor}{capacitor} de \qty{470}{\micro F} é carregado a \qty{12}{V}: que carga ele guarda? Que
tensão poria \qty{2.0}{mC} em \qty{100}{\micro F}?
\end{exercise}

\begin{exercise}[$\star$]\label{exo:g12:rc-rl-circuits:2}
Distribua \qty{47}{pF}, \qty{2200}{\micro F} e \qty{10}{F} entre um \omterm{def:g11:color-light-sources:subtractive}{filtro} de fonte de
alimentação, um sintonizador de rádio e um supercapacitor; calcule a \omterm{def:g10:energy-conservation:energy}{energia} do último
a \qty{2.7}{V}.
\end{exercise}

\begin{exercise}[$\star$]\label{exo:g12:rc-rl-circuits:3}
A tensão num \omterm{def:g12:rc-rl-circuits:capacitor}{capacitor} de \qty{220}{\micro F} sobe de modo uniforme de $0$ a
\qty{5.0}{V} em \qty{10}{ms}: que corrente entra? E que corrente circula depois que a
tensão para de variar, e por quê?
\end{exercise}

\begin{exercise}[$\star$]\label{exo:g12:rc-rl-circuits:4}
Calcule $\tau$ para $R = \qty{10}{k\ohm}$ e $C = \qty{100}{\micro F}$, e depois para
$L = \qty{0.10}{H}$ e $R = \qty{50}{\ohm}$. Confira, \omterm{def:g10:orders-of-magnitude:unit}{unidade} por \omterm{def:g10:orders-of-magnitude:unit}{unidade}, que as duas
combinações dão segundos.
\end{exercise}

\begin{exercise}[$\star$]\label{exo:g12:rc-rl-circuits:5}
Um \omterm{def:g12:rc-rl-circuits:capacitor}{capacitor} carrega rumo a \qty{6.0}{V} e passa por \qty{3.8}{V} em
$t = \qty{2.0}{s}$. Leia a \omterm{def:g12:rc-rl-circuits:tau}{constante de tempo}; quando a carga está $95\%$ concluída?
\end{exercise}

\begin{exercise}[$\star\star$]\label{exo:g12:rc-rl-circuits:6}
Um \omterm{thm:g12:rc-rl-circuits:charging}{circuito RC}: $\mathcal{E} = \qty{9.0}{V}$, $R = \qty{4.7}{k\ohm}$,
$C = \qty{220}{\micro F}$. Calcule $\tau$, $u(\tau)$, $u(3\tau)$ e a corrente no
instante em que a chave fecha.
\end{exercise}

\begin{exercise}[$\star\star$]\label{exo:g12:rc-rl-circuits:7}
Verifique por substituição que $u(t) = \mathcal{E}(1 - e^{-t/RC})$ satisfaz
$RC\,du/dt + u = \mathcal{E}$; o que o valor dela em $t = 0$ e o limite para $t$ grande
codificam fisicamente?
\end{exercise}

\begin{exercise}[$\star\star$]\label{exo:g12:rc-rl-circuits:8}
Um \omterm{def:g12:rc-rl-circuits:capacitor}{capacitor} de \qty{100}{\micro F} carregado a \qty{12}{V} se descarrega através de
\qty{47}{k\ohm}: calcule $\tau$, $u(\tau)$ e o tempo para $u$ cair abaixo de
\qty{0.60}{V}.
\end{exercise}

\begin{exercise}[$\star\star$]\label{exo:g12:rc-rl-circuits:9}
A corrente numa \omterm{def:g11:magnetic-fields:solenoid}{bobina} de \qty{0.50}{H} sobe de modo uniforme de $0$ a \qty{2.0}{A} em
\qty{40}{ms}: que tensão aparece nos terminais dela? Com $u = L\,di/dt$, por que
\emph{abrir} uma chave sobre uma \omterm{def:g11:magnetic-fields:solenoid}{bobina} dá faísca, e fechá-la não?
\end{exercise}

\begin{exercise}[$\star\star$]\label{exo:g12:rc-rl-circuits:10}
O relé do \cref{ex:g12:rc-rl-circuits:rl}: verifique por substituição que
$i(t) = (\mathcal{E}/R)(1 - e^{-Rt/L})$ satisfaz $L\,di/dt + Ri = \mathcal{E}$ e depois
calcule $i(\tau)$ e $i(3\tau)$.
\end{exercise}

\begin{exercise}[$\star\star$]\label{exo:g12:rc-rl-circuits:11}
Mostre que $u'(0) = \mathcal{E}/\tau$ e deduza que a tangente na origem encontra a
assíntota em $t = \tau$. Numa curva registrada que dá $\tau = \qty{0.50}{s}$, com
$R = \qty{2.0}{k\ohm}$: encontre $C$.
\end{exercise}

\begin{exercise}[$\star\star\star$]\label{exo:g12:rc-rl-circuits:12}
Ordene pela \omterm{def:g10:energy-conservation:energy}{energia} armazenada, calculando cada uma: um supercapacitor de \qty{3000}{F}
a \qty{2.7}{V}; um \omterm{def:g12:rc-rl-circuits:capacitor}{capacitor} de flash, \qty{800}{\micro F} a \qty{300}{V}; uma \omterm{def:g11:magnetic-fields:solenoid}{bobina}
de \qty{10}{mH} a \qty{10}{A}; uma pilha AA ($E \approx \qty{1.3e4}{J}$). O que os três
primeiros oferecem que a pilha não consegue?
\end{exercise}

\begin{exercise}[$\star\star\star$]\label{exo:g12:rc-rl-circuits:13}
Um temporizador de escada mantém a luz acesa enquanto $u < 0.63\,\mathcal{E}$ no
\omterm{def:g12:rc-rl-circuits:capacitor}{capacitor} que carrega, apagando em $t = \tau$. Projete a luz para durar \qty{5.0}{s}
usando $C = \qty{100}{\micro F}$: qual é o $R$? Para um limiar de $0.50\,\mathcal{E}$,
mostre que o \omterm{def:g12:mechanical-waves:delay}{atraso} é $\tau \ln 2$ e calcule-o.
\end{exercise}

\begin{exercise}[$\star\star\star$]\label{exo:g12:rc-rl-circuits:14}
Um \omterm{def:g12:rc-rl-circuits:capacitor}{capacitor} de \qty{2200}{\micro F} filtra uma fonte: entre duas recargas, ele sozinho
alimenta uma carga que puxa \qty{0.50}{A} por \qty{10}{ms}. A partir de
$i = C\,du/dt$, calcule a queda de tensão e depois a \omterm{def:g12:rc-rl-circuits:capacitor}{capacitância} que a mantém abaixo de
\qty{0.50}{V}.
\end{exercise}

\begin{exercise}[$\star\star\star$]\label{exo:g12:rc-rl-circuits:15}
Uma \omterm{def:g11:magnetic-fields:solenoid}{bobina} de ignição, $L = \qty{0.80}{H}$, conduz \qty{2.0}{A} quando o platinado a
corta em cerca de \qty{1.0}{ms}. Estime a tensão na \omterm{def:g11:magnetic-fields:solenoid}{bobina} durante o corte e a \omterm{def:g10:energy-conservation:energy}{energia}
da faísca; por que uma interrupção lenta não dá faísca alguma?
\end{exercise}

\section{Problema: o flash da câmera}

\begin{problem}[{Problema de fim de semana --- o flash da câmera: dois segundos de
zunido, um milissegundo de glória e o primo de parede que faz corações voltarem a bater
--- uma só exponencial governa todos eles}]
\label{pb:g12:rc-rl-circuits:1}
Uma \omterm{def:g10:orders-of-magnitude:unit}{unidade} de flash compacta armazena \omterm{def:g10:energy-conservation:energy}{energia} num \omterm{def:g12:rc-rl-circuits:capacitor}{capacitor} $C = \qty{800}{\micro F}$,
carregado através de $R = \qty{1.0}{k\ohm}$ a partir de um \omterm{def:g10:energy-conservation:transfer}{conversor} modelado como fonte
ideal $\mathcal{E} = \qty{300}{V}$; a luz de ``pronto'' acende a $95\%$ da tensão
final. Disparado, o \omterm{def:g12:rc-rl-circuits:capacitor}{capacitor} enxerga apenas o tubo de flash,
$R_f = \qty{0.50}{\ohm}$. Dados: uma pilha AA entrega \qty{1.5}{V} e \qty{2.4}{A.h};
$e = \qty{1.60e-19}{C}$.

\medskip
\noindent\textbf{Parte I --- O reservatório.}
\begin{enumerate}
  \item Calcule a carga no \omterm{def:g12:rc-rl-circuits:capacitor}{capacitor} a \qty{300}{V}.
  \item Quantas \omterm{def:g11:fundamental-interactions:elementary}{cargas elementares} são essas?
  \item Calcule a \omterm{def:g10:energy-conservation:energy}{energia} armazenada $E_C$.
  \item A pilha AA guarda $q = It \approx \qty{8.6e3}{C}$, ou seja, cerca de
        \qty{1.3e4}{J}. Quantas vezes a \omterm{def:g10:energy-conservation:energy}{energia} do \omterm{def:g12:rc-rl-circuits:capacitor}{capacitor} é isso?
  \item Então por que se dar ao \omterm{def:g11:work-of-force:work}{trabalho} do \omterm{def:g12:rc-rl-circuits:capacitor}{capacitor}? Em uma frase, usando a palavra
        ``\omterm{def:g11:work-of-force:power}{potência}''.
\end{enumerate}

\medskip
\noindent\textbf{Parte II --- A espera pelo verde.}
\begin{enumerate}[resume]
  \item Escreva a lei das malhas e transforme-a na equação diferencial de $u(t)$.
  \item Verifique por substituição que $u(t) = \mathcal{E}(1 - e^{-t/RC})$ a resolve,
        partindo de $u(0) = 0$.
  \item Calcule a \omterm{def:g12:rc-rl-circuits:tau}{constante de tempo} $\tau$.
  \item Depois de quanto tempo a luz de pronto fica verde, e em que tensão? Reconheça o
        zunido de dois segundos.
  \item Calcule a corrente inicial de carga; por que ela então morre?
\end{enumerate}

\medskip
\noindent\textbf{Parte III --- O milissegundo de glória.}
\begin{enumerate}[resume]
  \item Escreva a lei de descarga $u(t)$ através do tubo e calcule a nova \omterm{def:g12:rc-rl-circuits:tau}{constante de
        tempo} $\tau'$.
  \item Calcule a corrente de pico no tubo.
  \item Calcule a \omterm{def:g11:work-of-force:power}{potência} de pico; compare-a com uma chaleira elétrica de
        \qty{2.2}{kW}.
  \item Tomando o flash como terminado em $3\tau'$, dê a duração e a \omterm{def:g11:work-of-force:power}{potência} média
        dele.
  \item Forme a razão entre o tempo de carga e o tempo de flash: mesmo \omterm{def:g12:rc-rl-circuits:capacitor}{capacitor}, mesma
        carga --- o que mudou, e o que isso faz de um \omterm{def:g12:rc-rl-circuits:capacitor}{capacitor}: um tanque de \omterm{def:g10:energy-conservation:energy}{energia}
        ou uma alavanca de \omterm{def:g11:work-of-force:power}{potência}?
\end{enumerate}

\medskip
\noindent\textbf{Parte IV --- O primo na parede.}
Um desfibrilador precisa entregar $E_C = \qty{200}{J}$ a $U_0 = \qty{1500}{V}$; o tórax
e as pás apresentam cerca de \qty{50}{\ohm}.
\begin{enumerate}[resume]
  \item Calcule a \omterm{def:g12:rc-rl-circuits:capacitor}{capacitância} necessária.
  \item Calcule a carga armazenada.
  \item Calcule a \omterm{def:g12:rc-rl-circuits:tau}{constante de tempo} e a duração do choque ($3\tau$); compare com o
        flash.
  \item Calcule a corrente de pico e a \omterm{def:g11:work-of-force:power}{potência} de pico no paciente.
  \item O circuito de recarga precisa deixar o aparelho pronto ($95\%$) em
        \qty{6.0}{s}: encontre a \omterm{def:g11:circuits-and-power:resistance}{resistência} de carga necessária e a corrente inicial de
        carga --- os números que uma projetista encomendaria.
\end{enumerate}
\end{problem}
